\documentclass[11pt]{article}

% --------------------------------------------------
% Packages
% --------------------------------------------------
\usepackage[margin=1in]{geometry}
\usepackage{amsmath, amssymb, amsthm}
\usepackage{mathtools}
\usepackage{graphicx}
\usepackage{float}
\usepackage{enumitem}
\usepackage{hyperref}
\usepackage{xcolor}
\usepackage{tikz}
\usetikzlibrary{arrows.meta, positioning, calc}

% Pseudocode and algorithms
\usepackage{algorithm}
\usepackage{algpseudocode}

% Better fonts (optional)
\usepackage{lmodern}

% --------------------------------------------------
% Hyperref setup
% --------------------------------------------------
\hypersetup{
    colorlinks=true,
    linkcolor=blue,
    citecolor=blue,
    urlcolor=blue
}

% --------------------------------------------------
% Theorem-like environments
% --------------------------------------------------
\theoremstyle{definition}
\newtheorem{definition}{Definition}[section]
\newtheorem{theorem}{Theorem}[section]
\newtheorem{claim}{Claim}[section]
\newtheorem{remark}{Remark}[section]

\theoremstyle{proof}
% Proof environment is built-in via amsthm

% --------------------------------------------------
% Custom commands
% --------------------------------------------------
\newcommand{\Gen}{\mathsf{Gen}}
\newcommand{\Enc}{\mathsf{Enc}}
\newcommand{\Dec}{\mathsf{Dec}}
\newcommand{\Sig}{\mathsf{Sig}}
\newcommand{\Ver}{\mathsf{Ver}}
\newcommand{\negl}{\mathsf{negl}}

% --------------------------------------------------
% Document
% --------------------------------------------------
\begin{document}

\title{Public Key Crypto notes - 07}
\author{By contributors}
\date{March 23. 2026}
\maketitle

% --------------------------------------------------
\section{Schnorr-signature recap}
% --------------------------------------------------

\begin{itemize}
	\item $\Gen$:
	\begin{itemize}
		\item Input: $1^n$
		\item Output: $pk = (p, g, q, h = g^x), sk = x$
	\end{itemize}
	\item $\Sig$:
	\begin{itemize}
		\item Input: $sk, m$
		\item Output: $(I, s)$, where:
		\begin{align}
			k &\in_R \mathbb{Z}_q, \\
			I &= g^k, \\
			r &= H(I, m), \\
			s &= rx + k
		\end{align}
	\end{itemize}
	\item $\Ver$:
	\begin{itemize}
		\item Input: $pk, m, (I, s)$
		\item Output: Compute $r = H(I || m)$ and check whether $I = g^s \cdot h^{-r}$
	\end{itemize}
\end{itemize}

This scheme is patented. Instead of this scheme, we mostly use DSA (made by NIST).

\section{Digital Signature Algorithm (DSA)}

\begin{itemize}
	\item Designed by NIST.
	\item Let $H: {0, 1}^* \rightarrow \mathbb{Z}_q, F: {0, 1}^* \rightarrow \mathbb{Z}_q$ be two functions. Consider the following three algorithms:
	\item $\Gen$:
	\begin{itemize}
		\item Input: $1^n$
		\item Output: Public key $pk = (p, g, q, h = g^x)$, Secret key: $sk = x$
	\end{itemize}
	\item $\Sig$:
	\begin{itemize}
		\item Input: $m, sk$
		\item Output: $(r, s)$
		\item $k \in_R \mathbb{Z}_q, r = F(g^k), s = k^{-1}(H(m) + xr) \pmod{q}$
	\end{itemize}
	\item $\Ver$:
	\begin{itemize}
		\item Input: $m, (r, s), pk$
		\item Output: $r \stackrel{?}{=} F(g^{H(m) \cdot s^{-1}} \cdot h^{r \cdot s^{-1}})$
	\end{itemize}
	\item Correctness:
	\begin{align}
		g^{H(m) \cdot s^{-1}} \cdot h^{r \cdot s^{-1}} &= \\
		g^{H(m) \cdot s^{-1} + x \cdot r \cdot s^{-1}} &= \\
		g^{s^{-1}(H(m) + x \cdot h)} &= \\
		g^{\frac{k}{H(m) + xr} \cdot (H(m) + xr)} &= g^k \\
		F(g^k) \checkmark
	\end{align}
\end{itemize}

\begin{theorem}[Security of the DSA scheme]
	If $H$ and $F$ are modelled by a random oracle and the discrete logarithm problem is hard, then DSA is secure.
\end{theorem}

\begin{remark}
	In practice, $F(x) = x \pmod q$, this means that the ”random oracle” is not really random.
\end{remark}

\begin{remark}
	If the same $k$ is used multiple times, then one can obtain $sk = x$. It is important to use a different $k$ for each iteration.
	
	\noindent Indeed:
	
	\begin{align}
		s_1 &= k^{-1} (H(m_1) + x + r) \\
		s_2 &= k^{-1} (H(m_2) + x + r) \\
		s_1 - s_2 &= k^{-1} (H(m_1) - H(m_2)) \rightarrow k \text{ can be computed}\\
		\rightarrow x &\text{ can also be computed}.
	\end{align}
	
	\noindent Note: Playstation 3: All $k$ used for updates were the same.\footnote{\url{https://deeprnd.medium.com/decoding-the-playstation-3-hack-unraveling-the-ecdsa-random-generator-flaw-e9074a51b831}}
\end{remark}

\subsection{Parameter selection}

\begin{itemize}
	\item $p$ is large, $\approx 3000-4000$ bits,
	\item $q$ is medium size (size $\approx 256$ bits),
	\item $\rightarrow$ signature size: $2 \times 256$ bits $\approx500$ bits.
\end{itemize}

\section{Further digital signatures using the discrete log problem}

General form of signatures:

\begin{itemize}
	\item Given $a, b, c: ak = b + cx \pmod q$
	\item Verification: $g^{ak} \stackrel{?}{=} g^b \cdot g^{xc}$
	\item Examples:
	\begin{enumerate}
		\item Schnorr:
		\begin{itemize}
			\item $a = -1$
			\item $b = s$
			\item $c = -r$
		\end{itemize}
		\item DSA:
		\begin{itemize}
			\item $a = s$
			\item $b = H(m)$
			\item $c = (q^k \pmod{p}) \pmod{q}$
		\end{itemize}
		\item ElGamal:
		\begin{itemize}
			\item $a = s$
			\item $b = h(m)$
			\item $c = -r \pmod{p-1}$
		\end{itemize}
	\end{enumerate}
\end{itemize}

\section{Exotic protocols}

\subsection{Blind signature}

\begin{itemize}
	\item Given: signer $S$ and user $U$.
	\item $U$ has a message $m$ and wants $S$ to sign it without telling $m$ to $S$.
\end{itemize}

\subsection{Blind RSA signature}

\begin{itemize}
	\item The signer has plain RSA keys, where $pk = (N, e), sk = d$.
	\item $U$ has message $m \in \mathbb{Z}_N^*$, blinding parameter $r \in_R \mathbb{Z}_N^*$, and sends $m' = m \cdot r^e \cdot N$ to $S$.
	\item $S$ signs $m'$, and sends $s' = m'^d \pmod{N}$ to $U$.
	\item So $s = s' \cdot r^{-1} \pmod{N} = m^d (r^e)^d r^{1} = m^d$
\end{itemize}

\begin{remark}
	The value $r$ is the blinding parameter. If it is chosen randomly, then $m'$ cannot leak any information about $m$.
\end{remark}

\subsection{Blind Schnorr signature}

\begin{itemize}
	\item $S$ chooses $k \in_R \mathbb{Z}_q$, sends $I=g^k$ to $U$.
	\item $U$ chooses blinding parameters $\alpha, \beta \in_R \mathbb{Z}_q$.
	\item $U$ calculates:
	\begin{align}
		I' &= I \cdot g^\alpha \cdot h^\beta \\
		r' &= H(I' || m) \\
		r &= r' + \beta
	\end{align}
	and sends $r$ to $S$.
	\item $S$ calculates $s = rx + k$ and sends it to $U$.
	\item $U$ calculates $s' = s + \alpha$
	\item The signature is $sig = (I', s')$
\end{itemize}

Correctness:

\begin{align}
	g^s &\stackrel{?}{=} I \cdot h^r = \\
	g^{s'-\alpha} &\stackrel{?}{=} I' \cdot g^{-\alpha} \cdot h^{-\beta} \cdot g^{r' + \beta} = \\
	g^{s'} &\stackrel{?}{=} I' \cdot h^{r'}
\end{align}

\subsection{Oblivious transfer}

There are two parties: sender $A$ and receiver $B$. $A$ wants to send a message (out of two) to $B$, such that:

\begin{itemize}
	\item $B$ can only receive one of the messages.
	\item $A$ should \emph{not} learn which message $B$ received.
	\item Formally:
	\begin{itemize}
		\item Two messages: $m_0, m_1$
		\item $B$ chooses $b \in_R {0, 1}$ and receives $m_b$
	\end{itemize}
\end{itemize}

\subsubsection{Protocol 1 (Even, Goldreich, Lempel)}

\begin{itemize}
	\item Based on RSA.
	\item $A$ has two messages: $m_0$, $m_1$ and generates RSA keys $pk = (N, e), sk = d$.
	\item $A$ chooses $x_0, x_1 \in_R \mathbb{Z}_N^*$ and sends them to $B$ along with $pk$.
	\item $B$ chooses $b \in_R {0, 1}$ and $k \in_R \mathbb{Z}_N^*$, and computes $v = (x_b + k^e) \pmod{N}$ and sends $v$ (also called a \emph{pre-envelope}) to $A$.
	\item $A$ computes:
	\begin{align}
		k_0 &= (v - x_0)^d, \\
		k_1 &= (v - x_1)^d, \\
		m'_0 &= m_0 + k_0, \\
		m'_1 &= m_1 + k_1, \\
	\end{align}
	and sends $m'_0, m'_1$ to $B$.
	\item $B$ computes $m_b = m_b' - k$
\end{itemize}

But $B$ cannot compute $k_{1-b} = (v - x_{1-b})^d$

\begin{remark}
	This is a one-out-of-two oblivious transfer, but there are 1-out-of-n or m-out-of-n oblivious transfer schemes out there also.
\end{remark}

\end{document}